\section{Experimental Setup}

Evaluation is dominated by one fact: entanglements are rare, contiguous events, and the 25 recordings differ in gait and affected leg. A window-level random split would put near-identical neighbouring windows from one recording on both sides of the partition, leaking information and inflating every score, so all protocols group by whole recording. We report two. The first is a fixed leakage-safe split of 16 training / 4 validation / 5 test recordings (the five test recordings are four positive and one negative); normalization statistics and the walking reference distribution for the severity index are fit on training rows only. The second, our headline generalization protocol, is leave-one-recording-out (LORO) cross-validation over the 15 positive recordings: each fold holds out one positive recording, trains on the rest, and we report the mean and standard deviation of the per-fold $F_1$. Its variance across folds is itself informative. Both use dense evaluation, sliding the 200-sample window with hop 1 so a decision is produced at every timestep, matching how the runtime consumes \texttt{/lowstate}; this also scores the hard event boundaries.

Metrics are standard, computed on the dense windows. Precision, recall, and their harmonic mean $F_1$ summarize the detection head; PR-AUC and ROC-AUC are threshold-free, with PR-AUC the more meaningful here because positives are rare. For attribution we use multi-label exact-match, which counts a window correct only when the predicted set of entangled legs equals the true set exactly. As a non-learned reference, a thigh-torque heuristic fires when a leg's resistive torque is high; we evaluate it under the identical dense, grouped protocol on the same windows, so the comparison in Table~\ref{tab:detection} isolates the value of the learned model.

Two further analyses use the same protocol. For the feature ablation we retrain from scratch on channel subsets (engineered only, raw only, raw+engineered without foot contacts, without the IMU, and the full set) and re-run LORO, so the resulting change in $F_1$ reflects each group's marginal contribution rather than a post-hoc importance estimate. For calibration we report the Brier score, the expected calibration error (ECE), and saturation (the fraction of raw probabilities pinned near $1.0$); temperature scaling \cite{guo2017calibration,niculescu2005predicting} is applied to de-saturate the output so a usable operating threshold exists, and we report the effect without assuming de-saturation implies better reliability. Latency is measured on the deployment ONNX runtime \cite{onnxruntime}, pinned to a single CPU thread, timing the per-window forward pass plus post-processing; this is a development-machine benchmark, not a measurement on the Go2's onboard computer.
